\section{Conclusion}
\label{sec:conclusion}
We asked whether an asymmetric-reward dueling Double DQN (E-DDQN) improves safety-prioritized cybersickness detection in VR. On the public UFFCSData gameplay dataset the answer is no. On logged data the detection problem is a contextual bandit, and the 5:1 asymmetric reward induces exactly the Bayes threshold of 5:1 cost-sensitive classification. Under a nested, session-grouped protocol fixed before testing, E-DDQN reached a mean cost over seeds of \EDDQNCostCostMean{} and ROC-AUC \EDDQNCostRocAucMean, did not outperform a standard DQN (cost \DQNCostCostMean, ROC-AUC \DQNCostRocAucMean) and had a higher cost than a cost-weighted random forest in every outer fold, marginally in one (cost \RFCostCostMean, ROC-AUC \RFCostRocAucMean; seed ensembles \PHEnsEDDQNCost{} vs.\ \PHEnsRFCost; Holm-adjusted $p=$\,\PHHolmRF{} by group bootstrap and \PHSfRFHolm{} by the more conservative exact sign-flip test). A re-tuning with two further random search draws, specified before it was run, found no draw in which E-DDQN had a significantly lower cost than the standard DQN, and each agent ranked seconds better in one draw; it also showed that E-DDQN's weak primary result reflects its search draw: across draws its within-fold ROC-AUC ranged from \SDEDDQNWfMin{} to \SDEDDQNWfMax{} and its cost from \SDEDDQNCostMin{} to \SDEDDQNCostMax, still above the random forest's as a point estimate. A post-hoc logistic regression on elapsed time alone reached ROC-AUC \PHTimeLRRocAucMean, no primary model ranked seconds significantly better, and no model combined high recall with a controlled false-alarm rate or was shown to improve on always alarming at the 5:1 ratio. Session grouping was essential to these conclusions: splits that share sessions between training and test give near-perfect scores on these data, whereas session-grouped evaluation of the same features gives ROC-AUC \LeakGrpRocAuc. Bounded telemetry perturbations changed ROC-AUC by at most \AttAucMaxAbsChange, which is not evidence of robustness: most selected configurations use exposure-type inputs, and an adaptive adversary still changed decisions at fixed thresholds. Progress on this problem is likely to require datasets with physiological or eye-tracking signals, verified participant identifiers and, for adversarial claims, manipulations applied to participants with measured effects, and evaluations that include grouped splits, trivial and exposure-time baselines, cost-weighted comparators and more than one random search draw per model.
